\SourceLeaf{432}%
\[
\begin{aligned}
Z&=\pm\frac{R}{\BareRoot{(PP+QQ+RR)}}\\
&=\pm\frac{1}{\BareRoot{(1+tt+uu)}},
\end{aligned}
\]
hence
\[
AC-BB=\frac{TV-UU}{(1+tt+uu)^{2}}.
\]

To \textls[100]{every} point of the curved surface there corresponds a particular point $L$ on the auxiliary sphere, by means of the normal drawn there and the radius of the auxiliary sphere parallel to it. Likewise, the points on the auxiliary sphere corresponding to all points of a \textls[100]{line} on the curved surface form a line corresponding to it, and a bounded figure on the curved surface will similarly correspond to a bounded figure on the auxiliary sphere. We shall take the area of the latter as the measure of the amplitude of the former. We shall either regard this area as a number, taking the square of the radius of the auxiliary sphere as the unit, or express it in degrees and so on, setting the area of a hemisphere $=360^{0}$.

Comparing the area on the curved surface with the corresponding amplitude leads to the concept that we call the measure of curvature of the surface. If the former is proportional to the latter, the curvature is called uniform, and the quotient obtained by dividing the amplitude by the area is called the measure of curvature. This is the case when the curved surface is a sphere. The measure of curvature is then a fraction whose numerator is $1$ and whose denominator is the square of the radius.

We shall take the measure of curvature to be positive if the boundaries run around the figures on the curved surface and on the auxiliary sphere in the same sense, and negative if they surround the figures in opposite senses. If that proportionality does not hold, the surface has nonuniform curvature. At each point there is a particular measure of curvature, obtained by comparing corresponding infinitesimal parts of the curved surface and the auxiliary sphere. Let $d\sigma$ be a surface element on the former and $d\Sigma$ the corresponding element on the auxiliary sphere. Then
\[
\frac{d\Sigma}{d\sigma}
\]
is the measure of curvature at that point.
